\documentclass[10pt]{beamer}
\usetheme{Boadilla}
\usecolortheme{seahorse}
\setbeamertemplate{navigation symbols}{}
\setbeamertemplate{itemize items}[circle]
\usepackage{amsmath,amssymb}
\usepackage{booktabs}
\newcommand{\code}[1]{\texttt{#1}}
\newcommand{\wrapop}{\operatorname{wrap}}
\newcommand{\E}{\mathbb{E}}
\newcommand{\RN}[1]{\textcolor{blue!50!black}{[RN23#1]}}
\newcommand{\proj}[1]{\textcolor{green!35!black}{[proj#1]}}
\newcommand{\srcline}[1]{\par\vfill\noindent{\tiny\textcolor{gray}{Source: #1}}}

\title[OTA Carrier Phase Sync]{Over-the-Air Carrier Phase Synchronization\\ for Distributed Arrays}
\subtitle{A lecture: foundations, the reference paper, and this project}
\author[C.~Stevenson]{Chandler Stevenson}
\institute[]{\footnotesize Reference paper [RN23]:\\
M.~Rashid and J.~A.~Nanzer, ``Frequency and Phase Synchronization in
Distributed\\ Antenna Arrays Based on Consensus Averaging and Kalman
Filtering,''\\
\emph{IEEE Trans.\ Wireless Commun.}, vol.~22, no.~4, pp.~2789--2802,
Apr.~2023.\\ DOI 10.1109/TWC.2022.3213788}
\date{July 29, 2026}

\begin{document}

\begin{frame}
\titlepage
\end{frame}

\begin{frame}{The problem in one paragraph}
Two radio stations, A and B, each generate a 915\,MHz carrier from their
own crystal oscillator. We want their two carriers to hold a fixed,
known phase relationship --- within a few degrees --- so that their
signals add constructively at a distant point, or so that what the two
stations \emph{receive} can be combined coherently.
\vspace{2ex}
\begin{itemize}
\item The two crystals disagree, and their disagreement drifts randomly.
\item The only connection between the stations is radio itself.
\item So station B must continuously measure its offset from A
\emph{over the air} and steer its own electronics to cancel it.
\end{itemize}
\vspace{2ex}
Everything in this lecture is about how well that can be done, what
fundamentally limits it, and how the reference paper and this project
each attack it.
\end{frame}

\begin{frame}{Roadmap and attribution}
\begin{enumerate}
\item \textbf{Foundations} (21 slides) --- built up slowly; classical sources named per slide
\item \textbf{The reference paper} \RN{} (14 slides) --- its model, algorithms, derivation, and results taught with its own equation/figure numbers
\item \textbf{This project} \proj{} (19 slides) --- tools and method first, then every number names the command that regenerates it
\end{enumerate}
\vspace{1ex}
\begin{block}{Tags}
\RN{, Eq.~27} $=$ the paper, that equation.\quad
\proj{: \code{--model compare}} $=$ measured here, by that command.\quad
Named citations $=$ classical results.
\end{block}
\srcline{Companion with full derivations: \code{teaching\_document.pdf} (same numbers, same order).}
\end{frame}

%=====================================================
\section{Foundations}

\begin{frame}{1.0\quad Vocabulary for the whole lecture}
\footnotesize
\begin{description}
\item[Carrier / phase] A radio transmits a pure high-frequency sine wave
(here 915 million cycles per second) and imprints information on it.
\emph{Phase} is where in its cycle the wave is right now --- an angle,
one full cycle $=2\pi$ radians $=360^\circ$.
\item[Units] 1\,rad $=57.3^\circ$; \ 1\,mrad $=0.057^\circ$ (so
``28\,mrad'' means ``$1.6^\circ$''). \ Decibels compare powers on a log
scale: 20\,dB $=100\times$; ``20\,dB SNR'' means the signal is
$100\times$ stronger than the noise.
\item[Baseband / IQ samples] The receiver strips off the fast carrier
and keeps, at each sampling instant, one \emph{complex number} whose
angle is the received phase and whose magnitude is the strength. All
digital processing operates on these ``IQ samples.''
\item[Channel] Everything between the antennas: the travel delay plus
reflections off objects (\emph{multipath}). At one frequency its net
effect is a single complex multiplier --- a gain and a phase shift
$\phi_c$. \textbf{AWGN}: the electronic hiss every receiver adds,
modeled as Gaussian noise.
\item[Pilot / correlation] A \emph{pilot} is a pre-agreed known
waveform used for measurement. \emph{Correlating} means sliding the
known pattern along the received samples, multiplying and summing: the
result peaks when aligned, and the peak's \emph{angle} is the received
phase.
\end{description}
\srcline{Standard definitions; kept on one slide so later slides can use the words freely.}
\end{frame}

\begin{frame}{1.1\quad Two waves adding: where the requirement comes from}
A wave of one fixed frequency is completely described by an arrow (a
\emph{phasor}): length $=$ amplitude, angle $=$ phase; written as a
complex number the arrow is $e^{j\phi}$. Two waves arriving at the same
point add as arrows, so two equal-strength carriers give a sum that
depends only on the angle between them:
\[
|1+e^{j\Delta\phi}|^2 \;=\; (1{+}\cos\Delta\phi)^2+\sin^2\Delta\phi
\;=\; 2+2\cos\Delta\phi \;=\; 4\cos^2\!\big(\tfrac{\Delta\phi}{2}\big)
\]
\begin{itemize}
\item $\Delta\phi=0$: the phasors line up; power is $4\times$ one
transmitter. Two transmitters, four times the power --- the extra factor
is the point of coherence. With $N$ elements: $N^2$.
\item $\Delta\phi=\pi$: the phasors oppose; the sum is \emph{zero}.
Being wrong by half a carrier cycle is worse than transmitting alone.
\end{itemize}
\vspace{1ex}
We will use $G=\cos^2(\Delta\phi/2)$, the fraction of ideal two-element
gain, as the scorecard for every scheme in this lecture.
\srcline{Elementary phasor arithmetic.}
\end{frame}

\begin{frame}{1.2\quad How much phase error is tolerable?}
Let each element's phase error $\phi_n$ be random, zero-mean, Gaussian
with standard deviation $\sigma$. The mean of a unit phasor with jittered
angle is shorter than 1:
\[
\E[e^{j\phi_n}]=e^{-\sigma^2/2}
\quad\Longrightarrow\quad
\text{mean coherent power}\;\approx\;N^2\,e^{-\sigma^2}.
\]
Why the average arrow is shorter: wobbling the angle moves the arrow
mostly \emph{sideways}, and sideways excursions average to zero. But the
forward component is $\cos\phi_n$, which is $\le1$ for every wobble, so
its average only loses: $\E[\cos\phi_n]\approx1-\sigma^2/2$ for small
errors --- which is $e^{-\sigma^2/2}$. Each element still radiates full
power; the \emph{resultant} shrinks by that factor per element, hence
$e^{-\sigma^2}$ on power.
\vspace{1ex}
\begin{block}{The working threshold}
$\sigma=18^\circ=0.314$\,rad $\Rightarrow e^{-\sigma^2}=0.906$: with
$18^\circ$ RMS phase error an array keeps ${\ge}90\%$ of ideal gain.
This is the standard bar in the distributed-array literature.
\end{block}
For reference: $18^\circ=314$\,mrad. The loops later in this lecture
reach 23--82\,mrad --- comfortably inside.
\srcline{Threshold: Nanzer et al., IEEE TMTT 65(5), 2017 --- [RN23]'s reference [4], used throughout it.}
\end{frame}

\begin{frame}{1.3\quad Why this is hard: the scale of a carrier cycle}
\begin{itemize}
\item At 915\,MHz one carrier period is 1.09\,ns. A timing disagreement
of \textbf{1\,ns} between stations is $0.915$ of a cycle $=329^\circ$ of
phase --- more than a full tolerance budget, from one nanosecond.
\item Crystals are specified in parts per million. A \textbf{1\,ppm}
frequency error at 915\,MHz is 915\,Hz: the relative phase sweeps
through a \emph{full cycle every 1.09\,ms}.
\item So ``synchronize once'' is meaningless. The relative phase must be
measured and corrected \emph{continuously}, and the interesting question
becomes: what is the steady-state error of that loop?
\end{itemize}
\vspace{1ex}
Also note what coherent \emph{reception} needs: when two stations
receive the same signal from a passive source, the source's phase is
common to both and cancels when combining. What does not cancel is the
phase between the two stations' local oscillators. Station-to-station
sync is the whole requirement.
\srcline{Arithmetic from the carrier frequency; sensing observation: standard.}
\end{frame}

\begin{frame}{2.1\quad Modeling an oscillator}
An oscillator's deviation from its nominal frequency is described by two
numbers that evolve in time: a phase offset $\phi$ (rad) and a frequency
offset $\omega$ (rad/s). Between correction epochs spaced $T$ apart,
phase integrates frequency, and both get kicked by noise:
\[
\begin{bmatrix}\phi_{k+1}\\ \omega_{k+1}\end{bmatrix}
=\begin{bmatrix}1 & T\\ 0 & 1\end{bmatrix}
\begin{bmatrix}\phi_{k}\\ \omega_{k}\end{bmatrix}+\mathbf w_k .
\]
\begin{itemize}
\item The matrix says: new phase $=$ old phase $+$ $T\times$frequency;
frequency persists. This deterministic part is easy --- a controller can
predict and cancel it.
\item Everything difficult lives in $\mathbf w_k$: the random component
of real oscillator behavior. The next three slides unpack it.
\end{itemize}
\srcline{Standard clock model; used identically in [RN23] Eq.~28 and in this project.}
\end{frame}

\begin{frame}{2.2\quad The dominant noise: white FM $=$ a random walk in phase}
\small
``White FM'' means the oscillator's \emph{frequency} jitters
independently from sample to sample. Phase is the integral of frequency,
so phase takes a small independent random step each sample --- a random
walk. If each step has standard deviation $\sigma_{\mathrm{pn}}$ at
sample rate $f_s$:
\[
\operatorname{var}\bigl(\phi(t)-\phi(0)\bigr)
=\sigma_{\mathrm{pn}}^2 f_s\,t \;\equiv\; k\,t
\qquad\text{(variance grows linearly in time)} .
\]
\begin{alertblock}{Why this sets a floor no algorithm can beat}
A random walk has \emph{independent increments}: the steps taken after
your last measurement are statistically unrelated to everything you
observed before it. Whatever the walk accumulates between corrections
is therefore \textbf{unpredictable}, and a loop that corrects every $T$
seconds carries at least $kT$ of phase variance. More SNR, better
estimators, smarter filters --- none of it touches this term. Only a
quieter oscillator or a shorter $T$ does.
\end{alertblock}
\proj{} default: $\sigma_{\mathrm{pn}}=2\times10^{-4}$\,rad/sample at
1\,MS/s $\Rightarrow k=0.04$\,rad$^2$/s $\Rightarrow \sqrt{kT}=45$\,mrad
at $T=50$\,ms. Remember this number.
\srcline{Random-walk properties: standard. Parameter: \code{ota\_sync/sdr.py} defaults.}
\end{frame}

\begin{frame}{2.3\quad The other noise types}
\small
\begin{center}
\footnotesize
\begin{tabular}{llll}
\toprule
type & what it is & effect & ADEV\\
\midrule
white PM & i.i.d.\ jitter per phase sample & averages away as $1/L$ & $\tau^{-1}$\\
white FM & random-walk phase (prev.\ slide) & the $kT$ floor & $\tau^{-1/2}$\\
flicker FM & $1/f$ frequency noise & slow wander, no exact model & $\tau^{0}$\\
RW FM & random walk \emph{in frequency} & aging, temperature & $\tau^{+1/2}$\\
\bottomrule
\end{tabular}
\end{center}
\begin{itemize}
\item \textbf{White PM}: the phase \emph{reading} jitters but nothing
accumulates --- averaging many samples suppresses it. It matters only
inside individual measurements, never for long-term tracking.
\item \textbf{Flicker FM}: frequency noise with a $1/f$ spectrum ---
equal noise power per \emph{decade} of timescale, so it neither
averages away nor settles into a clean random walk. No exact
finite-state model exists; the standard surrogate is a sum of a few
first-order low-pass (AR(1)) processes with time constants spread
across the timescales of interest, whose summed spectra approximate
$1/f$.
\item \textbf{Random-walk FM}: the frequency itself wanders
(temperature, aging), so phase \emph{accelerates} away --- the reason
frequency must be tracked as its own state, not inferred from phase
alone.
\end{itemize}
\srcline{Taxonomy: IEEE Std 1139.}
\end{frame}

\begin{frame}{2.4\quad Grading oscillators: the Allan deviation}
\small
Ordinary variance fails for clocks: a drifting oscillator's frequency
has no settled mean, so ``the variance of $y$'' just grows with how
long you watch. The field's fix is the \textbf{Allan deviation}
$\sigma_y(\tau)$:
\begin{enumerate}
\item Take the fractional frequency $y=\dot\phi/(2\pi f_c)$
(dimensionless --- ``parts in $10^{11}$'' language).
\item Average it over consecutive windows of length $\tau$.
\item $\sigma_y(\tau)$ $=$ the typical difference between
\emph{successive} window averages (divided by $\sqrt2$).
\end{enumerate}
Differencing neighboring windows cancels slow drift, so the statistic
converges --- and each noise type leaves its own slope vs.\ $\tau$
(previous slide's table), making an ADEV plot a fingerprint of the
noise mix.
\begin{block}{How to read the number used throughout this lecture}
$\sigma_y(1\,\mathrm s)=3.5\times10^{-11}$ means: two successive
1-second frequency averages typically differ by 3.5 parts in $10^{11}$
--- at 915\,MHz, by 0.032\,Hz. Grades: \textbf{TCXO}
(temperature-compensated crystal, the few-dollar part in most radios)
sits at $\sigma_y\sim10^{-9}$--$10^{-10}$; \textbf{OCXO} (crystal kept
in a small oven) is 10--100$\times$ steadier but bigger, costlier,
power-hungry. \proj{} default sits between: realistic, not exotic.
\end{block}
\srcline{ADEV: IEEE Std 1139.}
\end{frame}

\begin{frame}{2.5\quad One crystal, three synchronized errors}
Inside a radio, the carrier synthesizer \emph{and} the ADC/DAC sample
clock are both multiplied up from the same reference crystal. So if the
reference runs fractionally fast by $y$, three things happen at once:
\[
\underbrace{\text{CFO}=f_c\,y}_{\text{carrier high}}\qquad
\underbrace{\text{SFO}=10^6 y\ \text{ppm}}_{\text{samples too fast}}\qquad
\underbrace{\text{timing drift}=y f_s T\ \text{samples per interval}}_{\text{frames arrive early/late}}
\]
At the \proj{} defaults, an initial CFO of 1500\,Hz at 915\,MHz means
$y=1.64\times10^{-6}$, hence 1.64\,ppm of SFO and 0.082 samples of frame
drift per 50\,ms --- \textbf{one} random process seen three ways, moving
together as the crystal wanders.
\vspace{1ex}
\begin{itemize}
\item A faithful simulation must derive all three from one process
(many models set them independently --- then they decorrelate, which
real hardware cannot do).
\item A digital correction (an NCO, introduced shortly) shifts the
\emph{baseband} phase/frequency; it does not touch the crystal. The
physical SFO and timing drift continue regardless of corrections.
\end{itemize}
\srcline{Shared-reference architecture: standard SDR design; modeled in \code{ota\_sync/sdr.py}.}
\end{frame}

\begin{frame}{3.1\quad Measuring a phase: the pilot and the bound}
Station A transmits a \emph{pilot}: a waveform B knows in advance. B
receives $L$ samples $r_n=e^{j\phi}s_n+w_n$ (known $s_n$, noise $w_n$)
and asks: what is $\phi$? Correlating $r$ against $s$ and taking the
angle is the maximum-likelihood answer.

How good can any estimator be? The Cram\'er--Rao lower bound (CRLB) is
the information-theoretic answer. Intuition: plot how probable the
observed data is as a function of the assumed $\phi$ --- a curve peaking
near the true value. The \emph{sharper} that peak, the more precisely
the data pins $\phi$ down; ``Fisher information'' is exactly that
sharpness (the curvature at the peak), and no unbiased method can have
variance smaller than its inverse:
\[
\operatorname{var}(\hat\phi)\;\ge\;\frac{1}{2\,\mathrm{SNR}\,L}.
\]
Each extra sample and each dB of SNR steepen the peak: twice the SNR or
twice the samples $\Rightarrow$ half the variance.
\begin{block}{At this project's settings}
20\,dB SNR ($\mathrm{SNR}=100$), $L=2\times2047$: bound $=1.1$\,mrad.
Phase can be measured \emph{very} precisely --- if noise were the only
problem.
\end{block}
\srcline{CRLB for tone parameters: Rife \& Boorstyn, IEEE Trans.\ Inf.\ Theory, 1974; also [RN23] Eq.~17.}
\end{frame}

\begin{frame}{3.2\quad Measuring frequency: phase slope over a baseline}
Frequency is the rate of phase change, so measure phase twice and
divide by the time between:
\[
\hat\omega=\frac{\hat\phi_2-\hat\phi_1}{T_{\mathrm{sep}}}
\qquad\Rightarrow\qquad
\operatorname{var}(\hat\omega)=\frac{1}{\mathrm{SNR}\,L_b\,T_{\mathrm{sep}}^2}.
\]
\begin{itemize}
\item The $T_{\mathrm{sep}}^2$ in the denominator is the useful lever: a
\emph{longer baseline} between the two phase looks improves frequency
quadratically. (Same reason a long telescope baseline improves angular
resolution.)
\item But the baseline cannot exceed the frame, and the phase difference
must stay unambiguous ($|\omega|T_{\mathrm{sep}}<\pi$) --- so receivers
use a \emph{coarse} estimate from short spacing first, then a
\emph{fine} one from long spacing. We will see exactly this two-stage
structure in the receiver.
\item At \proj{} settings ($L_b{=}2047$, $T_{\mathrm{sep}}{=}2.175$\,ms,
20\,dB): $\sigma_{\hat\omega}=0.16$\,Hz.
\end{itemize}
\srcline{Two-block frequency estimation: standard; variance form as implemented in \code{ota\_sync/sdr.py}.}
\end{frame}

\begin{frame}{3.3\quad The catch: your own clock jitters while you measure}
The CRLB assumed the only corruption is receiver noise. But the
receiver's oscillator runs its random walk \emph{during} the pilot, and
the phase estimate effectively averages that walk over the observation.
For a walk observed over samples $[a,\,a{+}n]$:
\[
\operatorname{var}(\overline W)\;\approx\;\sigma_{\mathrm{pn}}^2\,(a+n/3)
\]
(the walk arrives at the window already displaced --- the $a$ --- and
averaging a wandering path leaves $n/3$).
\begin{center}
\small
\begin{tabular}{lcc}
\toprule
per-measurement error & phase & frequency\\
\midrule
AWGN CRLB & 1.1\,mrad & 0.16\,Hz\\
own LO phase noise during pilot & \alert{8.3\,mrad} & \alert{0.68\,Hz}\\
\bottomrule
\end{tabular}
\end{center}
\begin{itemize}
\item Measurement here is \textbf{phase-noise limited, not SNR limited}:
more transmit power stops helping almost immediately.
\item A tracking filter must be told this (its $R$ must include the
phase-noise terms), or it will trust each measurement $\sim$7$\times$
too much and mis-weight everything.
\end{itemize}
\srcline{Walk-average variance and evaluation: this project, \code{teaching\_document.pdf} \S3.3.}
\end{frame}

\begin{frame}{3.4\quad What to transmit: Zadoff--Chu pilots \proj{}}
A good pilot needs two properties, and Zadoff--Chu (ZC) sequences,
$x[n]=e^{-j\pi u n(n+1)/L}$, deliver both:
\begin{enumerate}
\item \textbf{Constant amplitude} --- every sample has $|x[n]|=1$, so
the power amplifier runs at a steady level without distortion.
\item \textbf{Ideal autocorrelation} --- correlated against a shifted
copy of itself, a ZC gives zero at every nonzero (cyclic) shift. A
matched filter therefore produces \emph{one} sharp peak, giving precise
timing even when multipath smears the signal. A cyclic prefix (CP; the
sequence's tail copied in front) makes the channel's smearing look
cyclic, preserving the property.
\end{enumerate}
\vspace{1ex}
\proj{} frame (4606 samples $=$ 4.6\,ms at 1\,MS/s):
\[
\underbrace{\text{ZC-16}\times16}_{\text{STF: 256 samples}}
\;\rightarrow\;
\underbrace{[\text{CP }128+\text{ZC-2047}]\times2}_{\text{LTF: 4350 samples}}
\]
Short repeated blocks for \emph{finding} the packet; two long blocks for
\emph{precision}. Next slide: why that split.
\srcline{ZC: Chu, IEEE Trans.\ Inf.\ Theory, 1972. Note [RN23] transmits no waveform at all (its Eq.~15 observes oscillator states directly).}
\end{frame}

\begin{frame}{3.5\quad The receiver, stage by stage \proj{}}
\begin{enumerate}
\item \textbf{Detection + coarse CFO (from the STF).} Correlate the
signal against \emph{itself} delayed by one repetition (16 samples). A
repeated waveform matches itself regardless of channel or frequency
offset --- that is what makes detection robust. And a frequency offset
rotates each repetition relative to the last by $\omega\cdot16T_s$, so
the \emph{angle} of that self-correlation reads off the CFO. Unambiguous
while the rotation is under half a turn: $|f|<f_s/32=31.25$\,kHz. That
is the acquisition range.
\item \textbf{Fine timing.} Cross-correlate against the full known
preamble; the ZC property gives one sharp peak (sample-accurate).
\item \textbf{Fine CFO (from the two LTFs).} Phase difference across the
2175-sample baseline --- the long-baseline trick of slide 3.2, now that
coarse CFO has removed the ambiguity risk.
\item \textbf{Phase.} After removing the measured CFO, matched-filter
the LTFs; the correlation angle is the frame phase.
\end{enumerate}
Each stage needs the one before it: detect $\to$ time $\to$ frequency
$\to$ phase.
\srcline{Repeated-block detection/CFO: Schmidl \& Cox, IEEE Trans.\ Commun., 1997. Implementation: \code{ota\_sync/sdr.py}.}
\end{frame}

\begin{frame}{3.6\quad What the measurement actually contains}
Here is the central obstacle of the whole subject. The phase B measures
is
\[
m \;=\;
\underbrace{\phi_{\mathrm{tx}}-\phi_{\mathrm{rx}}}_{\theta:\ \text{oscillators, what we want}}
\;+\;
\underbrace{\phi_c}_{\text{channel: delay + multipath}}
\;+\; n .
\]
The wave accumulated phase traveling from A to B. The receiver sees one
number; $\theta$ and $\phi_c$ enter it \emph{identically}. No amount of
processing of one-way data can tell them apart.
\vspace{1ex}
\begin{alertblock}{Consequence}
A one-way loop that drives $m\to0$ does not set $\theta=0$; it sets
$\theta=-\phi_c$. The receiving oscillator \emph{absorbs the channel
phase} as a hidden bias. Measured at \proj{} defaults: $\phi_c=+2.92$
rad --- so one-way sync leaves the oscillators $2.92$\,rad apart while
every internal indicator reads ``locked.''
\end{alertblock}
Whether that bias matters depends on the application --- next slide.
\srcline{Identifiability framing and the measured bias: this project (\code{--model sdr}).}
\end{frame}

\begin{frame}{3.7\quad When the bias is harmless, and when it is fatal}
\textbf{Harmless: a user closes the loop (cell-free MIMO).} If a user
measures pilots from both stations, its channel estimate for station B
--- called CSI, \emph{channel-state information} --- contains B's
oscillator bias, indistinguishable from the channel to B and corrected
together with it. Any \emph{constant} per-station phase
drops out. Only the \emph{drift} between CSI updates matters, and the
sync loop controls exactly that. One-way sync suffices.
\vspace{1ex}
\textbf{Fatal: nothing closes the loop (passive coherent sensing).}
Beamforming toward a computed direction, or coherently combining echoes
from a target that cannot talk back: the carriers must be aligned
\emph{at the antennas themselves}. A hidden $2.92$\,rad offset simply
points the array somewhere else --- and nothing in the system can see
it.
\vspace{1ex}
\begin{block}{So the question becomes}
Can the stations measure $\theta$ \emph{without} the channel phase?
Yes --- by using the channel twice.
\end{block}
\srcline{Application taxonomy: standard; quantified for both cases later in this deck.}
\end{frame}

\begin{frame}{3.8\quad Two-way exchange: canceling the channel with itself}
Radio channels are \emph{reciprocal}: the phase accumulated from A to B
equals the phase from B to A over the same path at the same time. So
exchange pilots both ways within a channel coherence time:
\[
m_{\mathrm f}=\theta+\phi_c+n_1
\qquad
m_{\mathrm r}=-\theta+\phi_c+n_2
\]
($\theta$ flips sign with direction; $\phi_c$ does not). Then:
\[
\tfrac12(m_{\mathrm f}-m_{\mathrm r})=\theta+\tfrac{n_1-n_2}{2}
\qquad
\tfrac12(m_{\mathrm f}+m_{\mathrm r})=\phi_c+\tfrac{n_1+n_2}{2}
\]
The half-difference is the pure oscillator offset; the half-sum is the
pure channel phase; the two are uncorrelated, each with half the noise.
\vspace{1ex}
\begin{alertblock}{Two costs, both of which will bite later}
(1) Phases are only known modulo $2\pi$, and \emph{halving} a
modulo-$2\pi$ quantity leaves two possible answers, $\pi$ apart. One
branch must be chosen. \quad (2) Only the \emph{air} is reciprocal: each
node's transmit chain and receive chain are different electronics, and
that per-node TX/RX difference survives the subtraction --- it needs a
one-time loopback calibration.
\end{alertblock}
\srcline{Two-way transfer: classical (time-transfer lineage). Both caveats as stated: \code{teaching\_document.pdf} \S3.5.}
\end{frame}

\begin{frame}{4.1\quad Closing the loop: the NCO}
B cannot bend its crystal. What it \emph{can} do is rotate its digital
baseband: multiply samples by $e^{-j(\hat\phi+\hat\omega t)}$ --- a
\textbf{numerically controlled oscillator (NCO)}. Corrections are
digital, quantized (here: phase steps of $2\pi/2^{16}$, frequency steps
of 0.01\,Hz), and take effect only after the frame has been processed
--- one interval of \emph{latency} in this model.
\vspace{1ex}
The loop, once per interval $T$:
\[
\text{capture}\ \to\ \text{estimate }(m,\hat\omega)\ \to\
\text{filter}\ \to\ \text{command NCO (loads next interval)}
\]
\vspace{0.5ex}
Two questions remain, and they are the last two foundations:
\begin{enumerate}
\item How should noisy measurements be turned into corrections?
(the Kalman filter, next slide)
\item What does the loop's delay cost? (slide after)
\end{enumerate}
\srcline{NCO discipline: standard SDR practice; parameters: \code{ota\_sync/sdr.py}.}
\end{frame}

\begin{frame}{4.2\quad The Kalman filter, as a trust dial}
\small
Each interval the filter holds a \emph{prediction} of $[\theta,\omega]$
and receives a \emph{measurement}. It blends them:
\[
\text{estimate}=\text{prediction}+\kappa\,(\text{measurement}-\text{prediction})
\]
The gain $\kappa\in[0,1]$ is not tuned by hand; it is computed from two
declared variances:
\begin{itemize}
\item $Q$ (process noise): \emph{how fast does the truth wander between
measurements?} --- the oscillator physics of Part 2.
\item $R$ (measurement noise): \emph{how noisy is one measurement?} ---
the estimation limits of Part 3, including the phase-noise terms.
\end{itemize}
Large $Q$/small $R$: trust the meter ($\kappa\to1$). Small $Q$/large
$R$: trust memory ($\kappa\to0$). After a few steps $\kappa$ and the
error variance settle at a \textbf{steady state} --- the equilibrium
where information gained per measurement balances truth-wander per
interval. That steady-state uncertainty is a computable number, and it
is about to matter. (One word to keep: the filter's belief just after
absorbing a measurement is its \emph{posterior}; ``posterior
uncertainty'' means the spread of that belief.)

Phase wraps at $\pm\pi$, so the filter observes $[\cos m,\sin m]$
instead of $m$ (no wrap discontinuity) and re-linearizes within each
update --- an \emph{extended} KF.
\srcline{Kalman filtering: standard; iterated-EKF and Joseph-form details in \code{ota\_sync/core.py}.}
\end{frame}

\begin{frame}{4.2b\quad The trust dial in action: a toy example}
\footnotesize
One tracked quantity. The truth wanders by variance $Q=4$ per step
(steps of typical size 2); each measurement has variance $R=5$ (typical
error 2.24). Track the filter's own uncertainty $P$:
\[
\text{predict: } P^-\!=P+Q
\qquad
\text{gain: } \kappa=\frac{P^-}{P^-+R}
\qquad
\text{update: } P\leftarrow(1-\kappa)P^-
\]
\begin{center}
\small
\begin{tabular}{ccccc}
\toprule
round & $P^-$ (before meas.) & $\kappa$ & $P$ (after meas.) & error std\\
\midrule
1 & 4.00 & 0.44 & 2.22 & 1.49\\
2 & 6.22 & 0.55 & 2.77 & 1.66\\
3 & 6.77 & 0.58 & 2.88 & 1.70\\
$\infty$ & 6.90 & 0.58 & 2.90 & 1.70\\
\bottomrule
\end{tabular}
\end{center}
Three things to take away:
\begin{itemize}
\item The uncertainty \textbf{settles in a few rounds}, wherever it
starts: the equilibrium where the wander added per step ($Q$) balances
the information each measurement removes.
\item Settled error (1.70) beats any \emph{single} measurement (2.24)
--- memory helps --- but cannot reach zero: fresh wander keeps arriving.
\item This settled number is exactly the ``steady-state posterior
uncertainty'' that the latency term on the next slide integrates into
phase.
\end{itemize}
\srcline{Standard scalar Kalman recursion; numbers chosen for arithmetic clarity.}
\end{frame}

\begin{frame}{4.3\quad What latency costs}
The correction computed from frame $k$ loads at frame $k{+}1$ ---
$T=50$\,ms later. The controller is not naive: it knows its latency and
commands the \emph{predicted} state, $\hat\theta+\hat\omega T$. The
deterministic drift is cancelled. What survives?
\[
\theta_{\text{actual}}-\theta_{\text{commanded}}
=\underbrace{(\omega-\hat\omega)\,T}_{\text{frequency error, integrated}}
+\underbrace{\text{new walk}}_{\text{already counted}}
\]
The frequency \emph{estimate error} integrates into phase over the wait.
Its size is the filter's steady-state frequency uncertainty
$\sigma_{\hat\omega}$ --- computable from $Q_\omega$ and $R_\omega$ by
iterating the variance recursion to its fixed point:
\[
Q_\omega=0.79,\ R_\omega=1.03\ (\text{rad/s})^2
\;\Rightarrow\;
\sigma_{\hat\omega}=0.77\ \text{rad/s}
\;\Rightarrow\;
\sigma_{\hat\omega}T=38\ \text{mrad}.
\]
\begin{block}{The pre-registered prediction \proj{}}
Walk floor $\sqrt{kT}=45$\,mrad (slide 2.2) $\;+\;$ latency floor
38\,mrad $\;+\;$ tracking ${\sim}12$\,mrad. Neither of the first two
depends on SNR. Part 3 of this lecture measures exactly these numbers.
\end{block}
\srcline{Riccati fixed point: standard; this decomposition of a closed sync loop: this project (no counterpart in [RN23], whose updates are instantaneous).}
\end{frame}

%=====================================================
\section{The reference paper}

\begin{frame}{5.1\quad The reference paper: what it is and why it exists}
\begin{block}{Full citation}
M.~Rashid and J.~A.~Nanzer, ``Frequency and Phase Synchronization in
Distributed Antenna Arrays Based on Consensus Averaging and Kalman
Filtering,'' \emph{IEEE Transactions on Wireless Communications},
vol.~22, no.~4, pp.~2789--2802, April 2023. (arXiv:2201.08931; PDF in
this repository.)
\end{block}
The gap it fills, from its introduction \RN{, \S I}:
\begin{center}
\small
\begin{tabular}{lll}
\toprule
prior approach & how nodes align & why it falls short\\
\midrule
closed-loop & feedback bits from the destination & destination must talk back\\
primary--secondary & one master broadcasts a reference & dies with the master\\
decentralized & neighbors average & \textbf{frequency only}, no phase\\
\bottomrule
\end{tabular}
\end{center}
\vspace{1ex}
\textbf{This paper's goal:} no destination, no master --- every node
aligns \emph{frequency and phase} purely by averaging with its
neighbors, plus a \textbf{closed-form prediction} (Eq.~27) of the
residual error that remains.
\end{frame}

\begin{frame}{5.2\quad \RN{} The network picture}
\textbf{The cast:} $N$ radios ($N=5$ to 400 in their experiments), each
with its own free-running oscillator. Draw an edge between any two that
can hear each other: an undirected graph $\mathcal G$. No leader, no
destination --- the only goal is that all $N$ agree \emph{with each
other}.
\vspace{1ex}
\begin{description}
\item[{Connectivity $c\in(0,1]$}] the fraction of possible edges that
exist. $c=1$: everyone hears everyone. $c=0.05$: each node hears only a
few neighbors --- the realistic, difficult case most of the paper is
about.
\item[Rounds] once per interval $T$, every node broadcasts, listens,
and updates --- all in lockstep. (This is why the paper assumes time
synchronization is already solved.)
\item[TDMA] neighbors take turns transmitting, so with $D$ neighbors
each one is heard for only $L/D$ samples --- less data per estimate,
degrading the measurement errors below \RN{, Eqs.~18--21}.
\end{description}
\srcline{[RN23] \S II--III; their simulation parameters: $f_c=1$\,GHz, $f_s=10$\,MHz, $T=0.1$\,ms, SNR $=0$\,dB.}
\end{frame}

\begin{frame}{5.3\quad \RN{} What each node measures}
\small
Each node broadcasts its oscillator's \emph{bare tone}; each neighbor
estimates that tone's frequency and phase from $L$ noisy samples:
\[
x_n(t)=e^{j(2\pi f_n t+\theta_n)}+\text{noise}
\qquad\text{\RN{, Eq.~15}}
\]
The paper takes the estimates to be exactly as good as theory allows
--- the CRLBs (slide 3.1's bound, plus its frequency analog):
\[
\sigma^m_\theta=\sqrt{\frac{2}{L\,\mathrm{SNR}}}
\qquad
\sigma^m_f=\sqrt{\frac{6}{(2\pi)^2L^3\,\mathrm{SNR}}}\cdot f_s
\qquad\text{\RN{, Eqs.~16--17}}
\]
(Why $L^3$ for frequency: $L$ samples to average the noise, times a
window $L$ long over which a \emph{rate} is measured --- slide 3.2's
baseline effect, squared.)
\begin{alertblock}{The model's deepest assumption --- look at Eq.~15 again}
The tone arrives \emph{with no channel}: no propagation delay, no
multipath. A real receiver would measure $\theta_n+\phi_c$ and could
not split them (slide 3.6); here $\phi_c$ simply does not exist ---
bundled into a constant ``assumed corrected'' by some other mechanism.
\end{alertblock}
\vspace{-2pt}
Part 3 measures what happens when this assumption meets a real channel.
\srcline{[RN23] \S II-B.}
\end{frame}

\begin{frame}{5.4\quad \RN{} The error budget, part 1: oscillator drift}
\textbf{Question this slide answers:} how much does each oscillator
wander between two updates?
\vspace{1ex}
Their spec is an Allan-deviation curve (slide 2.4) for an ordinary
quartz TCXO --- white FM ($\beta_1$) plus random-walk FM ($\beta_2$),
no flicker:
\[
\sigma_f=f_c\sqrt{\beta_1/T+\beta_2T},\qquad
\beta_1=\beta_2=5\times10^{-19}
\qquad\text{\RN{, Eq.~2}}
\]
At their setup ($f_c=1$\,GHz, $T=0.1$\,ms) this gives
$\sigma_f\approx70.7$\,Hz: between two updates, a node's frequency
typically moves by 70.7\,Hz. That motion costs \emph{phase} in two
ways per interval \RN{, Eqs.~3--4, 13}:
\begin{center}
\small
\begin{tabular}{llc}
\toprule
term & what it is & size\\
\midrule
drift phase\ $2\pi\sigma_fT$ & phase the new offset sweeps before the next update & $2.55^\circ$\\
intra-interval\ $\pi\sigma_fT$ & the frequency keeps moving \emph{during} the interval & $1.27^\circ$\\
\bottomrule
\end{tabular}
\end{center}
\srcline{[RN23] \S II-B; numerical evaluation: this deck.}
\end{frame}

\begin{frame}{5.5\quad \RN{} The error budget, part 2: the full ranking}
\small
Two more error sources complete the list:
\begin{itemize}
\item \textbf{Phase jitter}: the spec sheet's total phase-noise power,
$A=-53.46$\,dBc, converts to $\sigma_\theta=2.7$\,mrad $=0.15^\circ$
--- \emph{redrawn fresh each interval}, no memory \RN{, Eq.~5}
(contrast slide 2.2's continuous walk).
\item \textbf{Measurement errors}: slide 5.3's CRLBs at their numbers
($L=Tf_s=1000$ samples, SNR $=0$\,dB): $\sigma^m_f=123$\,Hz,
$\sigma^m_\theta=44.7$\,mrad.
\end{itemize}
\begin{block}{All five per-interval error terms, ranked}
\centering
\begin{tabular}{llc}
1. & frequency \emph{measurement}, integrated: $2\pi\sigma^m_fT$ & $4.44^\circ$\\
2. & phase \emph{measurement}: $\sigma^m_\theta$ & $2.56^\circ$\\
3. & oscillator drift: $2\pi\sigma_fT$ & $2.55^\circ$\\
4. & intra-interval drift: $\pi\sigma_fT$ & $1.27^\circ$\\
5. & jitter: $\sigma_\theta$ & $0.15^\circ$\\
\end{tabular}
\end{block}
The ranking explains the paper's agenda: the biggest term is
\emph{measurement} error --- only 1000 noisy samples per look. And a
measurement-dominated loop is exactly where memory (filtering) helps,
which is why the paper's second algorithm exists (slide 5.11).
\srcline{[RN23] Eqs.~5, 16--17; numerical evaluation: this deck.}
\end{frame}

\begin{frame}{5.6\quad \RN{} DFPC: the algorithm is one averaging step, repeated}
\small
One round of \emph{decentralized frequency and phase consensus}
(their Algorithm 1) --- three steps:
\begin{enumerate}
\item \textbf{Drift.} Every oscillator degrades: frequency picks up
$\mathcal N(0,\sigma_f^2)$; phase picks up intra-interval drift and
jitter.
\item \textbf{Measure.} Every node broadcasts its tone; every neighbor
estimates it, with errors $\mathcal N(0,(\sigma^m_f)^2)$ and
$\mathcal N(0,(\sigma^m_\theta)^2)$.
\item \textbf{Average.} Every node \emph{replaces} its own frequency
and phase by a weighted average of what it heard, plus its own value:
\[
\mathbf f(k)=\mathbf W\,\mathbf f(k{-}1),\qquad
\boldsymbol\theta(k)=\mathbf W\,\boldsymbol\theta(k{-}1)
\qquad\text{\RN{, Eq.~8}}
\]
(The matrix is bookkeeping: row $n$ is nonzero only at node $n$'s
neighbors, so everything is computed locally. No leader.)
\end{enumerate}
\vspace{0.5ex}
Two things step 3 quietly builds in:
\begin{itemize}
\item \textbf{No data link anywhere}: the broadcast waveform itself is
the entire message.
\item \textbf{No actuator}: the averaged value simply \emph{is} the
node's new state --- no NCO, no latency (slide 4.3's term absent by
construction).
\end{itemize}
\srcline{[RN23] Algorithm 1, \S III.}
\end{frame}

\begin{frame}{5.7\quad \RN{} The mixing matrix, built by hand}
The weights are Metropolis--Hastings: for an edge $(i,j)$,
$w_{ij}=1/(1+\max\{n_i,n_j\})$ where $n_i$ is node $i$'s degree; the
self-weight absorbs the remainder \RN{, Eq.~9}.
\vspace{1ex}
\textbf{Worked example} --- three nodes in a line, $A\!-\!B\!-\!C$
(degrees 1, 2, 1):
\[
w_{AB}=w_{BC}=\frac{1}{1+\max(1,2)}=\tfrac13
\qquad\Rightarrow\qquad
\mathbf W=\begin{bmatrix}\tfrac23&\tfrac13&0\\[2pt]
\tfrac13&\tfrac13&\tfrac13\\[2pt]0&\tfrac13&\tfrac23\end{bmatrix}
\]
Check the two properties that make consensus work:
\begin{itemize}
\item \textbf{Rows sum to 1}: each update is a proper average --- no
node can run away.
\item \textbf{Columns sum to 1 too} (doubly stochastic): the
\emph{network average} is preserved by every round. So if the values
converge to a common number, that number must be the average of where
they started.
\end{itemize}
The $\max$ in the denominator is the safety margin: a weight small
enough for the busier endpoint is safe for both.
\srcline{[RN23] Eq.~9; worked example: this deck.}
\end{frame}

\begin{frame}{5.8\quad \RN{} Why and how fast consensus converges}
\small
Each round is a \textbf{blurring pass} over the network: the overall
average passes through untouched, while every \emph{pattern of
disagreement} fades.
\vspace{0.5ex}
\begin{itemize}
\item Exactly: $\mathbf W$'s largest eigenvalue is 1, for the
``everyone equal'' direction; every disagreement pattern shrinks by at
least a factor $\lambda_2$ (the second-largest eigenvalue) per round.
Repeated rounds therefore drive every state to the initial average:
\[
\mathbf W^{I}\ \longrightarrow\ \tfrac1N\mathbf{11}^\top .
\]
\item \textbf{Dense graph}: $\lambda_2$ small --- disagreement dies in
a few rounds. \textbf{Sparse graph}: $\lambda_2\to1$ --- information
must creep across the network, convergence is slow.
\end{itemize}
\vspace{0.5ex}
What their figures show:
\begin{itemize}
\item An initial frequency spread of $\pm100$\,ppm collapses within
tens of iterations at $c=0.2$ \RN{, Figs.~1--2}.
\item Larger $N$ at fixed $c$ converges \emph{faster}: each node has
$D=c(N{-}1)$ neighbors, so its local average is better.
\item Converged error falls with $N$ toward a floor \RN{, Fig.~3} ---
set by the noise injected each round, which is the next slide.
\end{itemize}
\srcline{[RN23] \S III-A, Figs.~1--3.}
\end{frame}

\begin{frame}{5.9\quad \RN{} Why the error does \emph{not} go to zero}
Mixing alone would converge exactly and stay. But steps 1--2 of
\emph{every} round inject fresh drift and fresh measurement error. The
network settles at an \textbf{equilibrium between two rates}:
\begin{center}
mixing \emph{removes} disagreement at $\sim(1-\lambda_2)$ per round\\[2pt]
vs.\quad injection \emph{adds} $\sigma_e$ of new error per round
\end{center}
Formally, the paper sums what survives of each past injection
($\sum_i\mathbf W^{i+1}\mathbf e_{I-i}$) and bounds the residual,
relative to the network average, by
$\sigma_e\sqrt{\sum_m\lambda_2^{2m}}$ \RN{, Eqs.~23--26}.
\begin{itemize}
\item \textbf{Dense}: old errors are mixed away quickly ---
equilibrium well below $\sigma_e$.
\item \textbf{Sparse} ($\lambda_2\to1$): each error lingers until the
next replaces it --- equilibrium saturates \emph{at the injection
level itself}.
\end{itemize}
\begin{block}{The punchline}
In the sparse (realistic) limit, the residual is simply the
root-sum-of-squares of the five per-round terms of slide 5.5 ---
assembled on the next slide as Eq.~27.
\end{block}
\srcline{[RN23] \S III-A.}
\end{frame}

\begin{frame}{5.10\quad \RN{} Eq.~27, assembled from the numbers we already have}
\[
\sigma_{\phi,\mathrm{total}}
=\sqrt{\underbrace{(2\pi\sigma_fT)^2}_{\text{drift}}
+\underbrace{(2\pi\sigma^m_fT)^2}_{\text{freq.\ est.}}
+\underbrace{(\pi T\sigma_f)^2}_{\text{intra}}
+\underbrace{(\sigma^m_\theta)^2}_{\text{phase est.}}
+\underbrace{\sigma_\theta^2}_{\text{jitter}}}
\qquad\text{\RN{, Eq.~27}}
\]
Every term was already computed on slides 5.4--5.5; plug them in:
\[
\sqrt{(2.55^\circ)^2+(4.44^\circ)^2+(1.27^\circ)^2+(2.56^\circ)^2+(0.15^\circ)^2}
=5.9^\circ
\]
\begin{itemize}
\item Well under the $18^\circ$/90\% bar of slide 1.2 --- consensus
alone suffices \emph{in this model}.
\item The biggest term is frequency estimation ($4.44^\circ$): the
motivation for filtering, next slide.
\item Compare with slide 4.3's closed-loop floor: Eq.~27 has \textbf{no
latency term} (instantaneous updates), \textbf{no channel term} (Eq.~15
has no channel), and \textbf{no noise memory} (jitter redrawn each
round). These absences are exactly where Part 3 will press.
\end{itemize}
\srcline{[RN23] Eq.~27; term-by-term evaluation: this deck (matches the $5.9^\circ$ used in Part 3's validation).}
\end{frame}

\begin{frame}{5.11\quad \RN{} KF-DFPC: filter first, then average}
\small
Slide 5.5's budget is dominated by \emph{measurement} error, and the
truth changes slowly between looks --- exactly when the trust dial of
slide 4.2 says: use memory. KF-DFPC changes one thing: each node
\textbf{Kalman-filters its own measurements before the network
averages}.
\begin{enumerate}
\item \textbf{Filter} (per node): two states $[f_n,\theta_n]$.
$\mathbf Q$ comes from the drift/jitter statistics of slides 5.4--5.5
(its off-diagonal term: a frequency kick also moves the phase it
integrates into) \RN{, Eq.~29}; $\mathbf R$ $=$ the CRLBs \RN{,
Eq.~31}; standard predict/update \RN{, Eqs.~33, 35--36}.
\item \textbf{Average} (as before): consensus on the \emph{posterior
means}; the covariances are mixed with squared weights, so each
node's filter knows how uncertain its new prior is \RN{, Eqs.~38--39}.
\end{enumerate}
\begin{alertblock}{The two costs of that one change}
(1) Posterior means \emph{and covariances} must now travel between
nodes: KF-DFPC needs an error-free \textbf{data side channel}, which
plain DFPC did not. (2) \textbf{No closed form} is derived for its
residual --- it is characterized only by simulation against Eq.~27.
\end{alertblock}
\srcline{[RN23] \S IV, Algorithm 2.}
\end{frame}

\begin{frame}{5.12\quad \RN{} Results, part 1: convergence and baselines}
\textbf{Convergence speed} (threshold: total phase error under
$1^\circ$):
\begin{center}
\small
\begin{tabular}{lcc}
\toprule
 & DFPC & KF-DFPC\\
\midrule
$N{=}100$, connectivity 0.05 & 16 iterations & \textbf{2}\\
$N{=}65$, connectivity 0.05 & 158 iterations & \textbf{4}\\
\bottomrule
\end{tabular}
\end{center}
Filtering shortens convergence dramatically on sparse graphs ---
smoothed values mix far more efficiently than noisy ones
\RN{, Figs.~4, 8}. Fewer rounds also means less energy spent
broadcasting.
\vspace{1ex}
\textbf{Baselines beaten} \RN{, Figs.~7--9} --- three existing
distributed estimators:
\begin{itemize}
\item \textbf{DLMS} (diffusion least-mean-squares): adaptive averaging
with no state model --- worst everywhere (it cannot exploit the slow
dynamics).
\item \textbf{DKF} (diffusion KF): shares estimates but not
covariances --- slow to converge.
\item \textbf{KCIF} (Kalman consensus information filter): shares
measurements and covariances --- competitive at small $N$/high SNR, but
plateaus (suboptimal gain).
\end{itemize}
KF-DFPC wins at low SNR and sparse connectivity --- the practical
corner.
\srcline{[RN23] \S V, Figs.~4, 7--9.}
\end{frame}

\begin{frame}{5.13\quad \RN{} Results, part 2: the update-interval sweep}
Their Fig.~10 sweeps $T$ from 0.1\,ms to 3\,s and finds a
\textbf{U-shape} for DFPC, minimum near $T=20$\,ms. The shape follows
from the term scalings we computed:
\begin{itemize}
\item \textbf{Drift terms grow with $T$}: more time between updates
means more phase swept out ($2\pi\sigma_fT$).
\item \textbf{Estimation terms shrink with $T$}: a longer interval means
more samples per look ($L=Tf_s$), so $\sigma^m_f,\sigma^m_\theta$ fall.
\item The optimum sits where the two cross --- at their parameters,
$T\approx20$\,ms.
\end{itemize}
KF-DFPC removes most of the left (estimation-dominated) branch:
filtering substitutes memory for samples, so short-$T$ operation stops
being punished. Consistently, its converged error is nearly
\textbf{SNR-independent} across $-30$ to $+120$\,dB \RN{, Fig.~11} ---
the filter compensates measurement quality with memory, at the cost of
a few more convergence iterations at low SNR \RN{, Fig.~9}.
\vspace{1ex}
Note the design language: this figure is a \emph{cadence trade-off} ---
the same axis this project's designs will push on, but without latency,
channel, or airtime in the model.
\srcline{[RN23] \S V, Figs.~9--11.}
\end{frame}

\begin{frame}{5.14\quad \RN{} The assumptions, and what each one hides}
\begin{enumerate}
\item \textbf{No channel in the measurement} (Eq.~15). Hides the
identifiability problem (slide 3.6) entirely --- over the air, every
measurement would contain $\phi_c$.
\item \textbf{Instantaneous updates.} Hides the latency term (slide
4.3) --- in the paper, a computed correction simply \emph{is} the new
state.
\item \textbf{Jitter redrawn each interval; no flicker.} Hides noise
memory: real intra-frame noise carries into the next interval.
\item \textbf{Statistics-level validation.} Its simulations draw
Gaussian errors from the assumed model (Algorithm 1) --- they exercise
the algebra, not the physics.
\item \textbf{Time sync assumed done; side channel ideal.}
\end{enumerate}
\vspace{1ex}
None of these are algebra errors --- they are scope choices, and
standard ones. Part 3's simulator exists precisely to measure what they
cost.
\srcline{Assessment: this project (\code{paper\_review\_rashid\_nanzer.pdf}), confirmed against the full text.}
\end{frame}

%=====================================================
\section{This project}

\begin{frame}{\proj{} How the project is built: tools, and why each one}
\begin{description}
\item[Python 3.11 + PyTorch] All signal processing and simulation is
tensor math over complex numbers; PyTorch supplies that, plus seeded
random-number generators and an optional GPU path. Everything runs in
\textbf{double precision} (\code{float64}/\code{complex128}) --- the
questions here live at the milliradian level, and single precision
($\sim$$10^{-7}$ relative error) would leave too little headroom under
long accumulations.
\item[NVIDIA Sionna 2 (\code{sionna-no-rt})] Supplies the channel
physics: the 3GPP TR~38.901 TDL multipath profiles, Jakes-correlated
tap evolution under mobility, and the blocks that apply them to
waveforms. Reason: the channel is the one part of the physics for which
a standards-grade, widely-validated implementation exists ---
hand-rolling it would add exactly the kind of silent error this project
exists to avoid.
\item[Determinism by construction] Every run is seeded (PyTorch global
seed, a separate stream for the radio's noise draws, Sionna's own
seed): identical command $\Rightarrow$ byte-identical output --- which
is what made the ``rerun and diff'' verification step meaningful.
\end{description}
\srcline{\code{requirements.txt}: \code{sionna-no-rt==2.0.1}; \code{pyproject.toml}: Python $\geq$ 3.11.}
\end{frame}

\begin{frame}{\proj{} How the code is organized}
\small
\begin{center}
\footnotesize
\begin{tabular}{ll}
\toprule
module & contents\\
\midrule
\code{ota\_sync/core.py} & oscillator model, iterated EKF (shared by every loop)\\
\code{ota\_sync/sdr.py} & config, frame, radio link + impairments, receiver,\\
 & one-way loop, oracle grading path\\
\code{ota\_sync/coherent.py} & two-way loop, $\pi$-calibration, CSI-gain evaluation\\
\code{ota\_sync/dfpc.py} & \RN{} reimplementations: statistics level \emph{and} OTA\\
\code{ota\_sync/microsync.py} & two-tier micro-pilot loop\\
\code{hybrid\_calibration/} & the 3-state hybrid (a separate package)\\
\code{simulation.py} & one CLI; \code{--model} selects the loop (7 models $+$ \code{compare})\\
\code{tests/} & 18 \code{pytest} regression tests\\
\bottomrule
\end{tabular}
\end{center}
\begin{itemize}
\item One design rule: every loop reuses the \emph{same} oscillator,
channel, and receiver blocks --- so when two schemes differ, the
difference is the scheme, not the plumbing. (This is what made the
82.8-vs-80.8 equivalence result meaningful.)
\item Each finding in this lecture is frozen as a test: oracle
consistency, the ablation floors, anti-phase capture, the calibration
flip, the Doppler-prior effect --- so later changes cannot silently
undo them.
\end{itemize}
\srcline{Repository layout; run \code{python -m pytest} from the repo root.}
\end{frame}

\begin{frame}{\proj{} Working method: how results were kept trustworthy}
The recurring pattern, applied to every claim in Part 3:
\begin{enumerate}
\item \textbf{Predict first.} Hand-compute the expected number from the
physics before running (the 45\,mrad walk and 38\,mrad latency floors
of slide 4.3 predate every simulation that confirmed them).
\item \textbf{Isolate by ablation.} Every physical effect has a switch
(\code{--no-rf-impairments}, \code{--correction-latency 0},
\code{--micro-pilots M}, \code{--anchor-every K}) so each term of a
budget can be measured alone.
\item \textbf{Grade against independent truth} (the oracle path) ---
never against the estimator's own inputs.
\item \textbf{Verify surprises independently.} The anti-phase capture
was confirmed by hand arithmetic, by a four-line radio-free model, and
by a seed sweep against the predicted capture condition --- three
routes that share no code with the simulator's loop.
\item \textbf{Check reproducibility literally}: rerun and \code{diff}
(deterministic seeding makes outputs byte-identical).
\item \textbf{Cross-check budgets}: e.g., unfiltered consensus's
measured 0.442\,Hz frequency noise predicts a 146\,mrad residual;
observed 153.
\end{enumerate}
\srcline{The same discipline explains the one thing simulation cannot supply: hardware validation (open item).}
\end{frame}

\begin{frame}{\proj{} The simulator: test the physics, not the algebra}
Everything from Part 1 is implemented at the waveform level and nothing
is handed to the receiver:
\begin{center}
\footnotesize
\begin{tabular}{llll}
\toprule
carrier / rate & 915\,MHz / 1\,MS/s & interval $T$ & 50\,ms\\
channel & 3GPP TDL-D + shadowing & noise floor & \textbf{fixed} (so SNR fades)\\
oscillators & all four noise types, continuous & clock & one crystal (slide 2.5)\\
TX/RX chain & PA, DAC, IQ imbal., AGC, DC, ADC & NCO & quantized, latency 1\\
\bottomrule
\end{tabular}
\end{center}
\begin{itemize}
\item ``TDL-D'': a standardized statistical multipath profile published
by the cellular-standards body 3GPP (report TR 38.901), implemented via
NVIDIA's Sionna library; ``shadowing'': slow strength fades from
obstructions. The other rows are the standard analog imperfections of
any real transceiver.
\item The receiver runs the full stage chain of slide 3.5 on raw IQ ---
detection can fail, timing can slip, estimates carry their real noise.
\item The intra-frame phase-noise walk \emph{continues} into the next
interval (folded into the oscillator state): noise has memory, unlike
\RN{}.
\end{itemize}
\srcline{\code{ota\_sync/} package; architecture: \code{simulation\_design.pdf}.}
\end{frame}

\begin{frame}{\proj{} Grading honestly: the oracle path}
\begin{alertblock}{A trap this project fell into, then fixed}
Compare the filter's estimate to the measurement that just updated it,
and the residual is only $(1-\kappa)\times$ the innovation (slide 4.2:
the estimate was \emph{pulled toward} that measurement). Both can be far
from the truth while agreeing with each other. This code initially
reported 0.105\,mrad tracking error that way.
\end{alertblock}
\vspace{1ex}
Fix: every capture is generated \emph{twice}. The \textbf{oracle twin}
shares everything deterministic --- same frame, same channel
realization, same timing, same PA/DAC distortion --- and omits
everything stochastic (noise, phase noise, receiver impairments).
Running the same estimator on the twin yields the true observable
phase; all errors are graded against that.
\vspace{1ex}
Ground-truth tracking RMSE: \textbf{12.0--12.5\,mrad} ---
${\sim}100\times$ the self-graded figure. (Method contrast: \RN{}
validates against draws from its own assumptions; here the claims are
checked against an independent physical model.)
\srcline{\code{ota\_sync/sdr.py} (oracle path); \code{--model sdr}, seed 0.}
\end{frame}

\begin{frame}{\proj{} The error budget, measured by ablation}
Slide 4.3 predicted the floor before any simulation. Now measure it,
isolating each term by switching physics off:
\begin{center}
\small
\begin{tabular}{lcc|l}
\toprule
configuration & tracking & residual & what the row isolates\\
\midrule
full defaults & 12.0\,mrad & 69.6\,mrad & all terms together\\
RF impairments off & 1.1\,mrad & 34.3\,mrad & latency alone (predicted 38)\\
zero latency & 12.3\,mrad & 12.3\,mrad & tracking alone\\
\bottomrule
\end{tabular}
\end{center}
Reading the table: with impairments off, the walk and phase-noise terms
vanish and the 34.3\,mrad left is the latency term --- against slide
4.3's hand-computed prediction of 38. The 1.1\,mrad tracking in that
row sits exactly on the AWGN CRLB (slide 3.1). Full defaults combine
$45\oplus38\oplus12$ plus flicker/shadowing $\to$ 69.6.
\begin{block}{Consequences}
The floor is set by \textbf{oscillator quality, correction cadence, and
latency} --- not pilot energy. And the latency term is absent from
\RN{, Eq.~27} and (per a 19-source verified search) from the published
literature through Dec.~2025.
\end{block}
\srcline{\code{--model sdr} with \code{--no-rf-impairments} / \code{--correction-latency 0}; search: \code{LITERATURE\_REVIEW.md}.}
\end{frame}

\begin{frame}{\proj{} Two-way sync: the bias is gone}
Implementing the exchange of slide 3.8 in the simulator: forward and
reverse frames traverse the same channel realization within the
interval, and the EKF consumes the half-difference phase and frequency
(each with half the one-way noise).
\begin{center}
\small
\begin{tabular}{lcc}
\toprule
 & one-way loop & two-way loop\\
\midrule
oscillator residual & $-2.92$\,rad (hidden bias) & \textbf{81.8\,mrad RMS}\\
two-station gain & near 0 open-loop & 99.83\%\\
airtime & 9.6\% & 19.1\%\\
\bottomrule
\end{tabular}
\end{center}
The carriers are now genuinely aligned \emph{at the antennas} --- the
capability passive sensing requires --- at double the pilot airtime.
\vspace{1ex}

But slide 3.8's caveat (1) is still owed: halving a wrapped quantity
leaves two candidate answers, $\pi$ apart, and a branch must be picked.
The next slide traces what a wrong pick does.
\srcline{\code{ota\_sync/coherent.py}, \code{--model twoway}, seed 0.}
\end{frame}

\begin{frame}{\proj{} A wrong branch is a stable illusion}
The loop resolves the two candidates by picking the one nearest its
current belief. Trace a wrong pick at acquisition, step by step:
\begin{enumerate}
\item \textbf{Acquisition.} True offset $\theta=1.94$\,rad; candidates
$\{1.94,\ -1.20\}$. With no prior belief the rule ``nearest zero''
picks $-1.20$ --- wrong by exactly $\pi$. The filter now believes
$\hat\theta=\theta-\pi$.
\item \textbf{Correction.} The loop drives the \emph{believed} phase to
zero. Since belief is $\pi$ off, the \emph{true} relative phase parks
at $\pi$: perfect anti-phase.
\item \textbf{Every later measurement confirms the belief.} The
half-difference again yields two candidates, $\{\approx0,\ \approx\pi\}$;
the one nearest the (near-zero) belief is chosen --- always the shifted
one. Innovations are tiny, covariance settles, lock indicators glow
green.
\item \textbf{Nothing internal can reveal it.} The measurement simply
cannot distinguish $\theta$ from $\theta-\pi$ --- that information does
not exist inside the loop. Only an \emph{external} observation (does
the pair combine constructively?) carries the missing bit.
\end{enumerate}
This is why the resolution (next slide) must be a one-time
\emph{combining} check, not more filtering.
\srcline{Mechanism verified by direct trace in \code{teaching\_document.pdf} \S11.}
\end{frame}

\begin{frame}{\proj{} The $\pi$ ambiguity: how it hid, how it is resolved}
\begin{alertblock}{Why every early run was (accidentally) fine}
Phase advance per interval $=2\pi\times1500\,\mathrm{Hz}\times
50\,\mathrm{ms}=75.000$ carrier cycles --- an \emph{exact integer}. So
the initial 1.2\,rad offset reappeared unchanged at every acquisition,
and since $1.2<\pi/2$, ``pick the candidate nearest zero'' happened to
be right every time. A later configuration used 5\,ms steps: 7.5 cycles
$\Rightarrow$ the acquisition phase landed at 1.94\,rad ($>\pi/2$)
$\Rightarrow$ wrong branch $\Rightarrow$ stable lock at $\pi$, 0.01\%
gain.
\end{alertblock}
\vspace{1ex}
Resolution (as a deployment would): after lock, one coarse combining
check --- constructive or destructive is \emph{one bit} of external
information --- and flip the NCO by $\pi$ if destructive.
\vspace{1ex}
One further subtlety, learned the hard way: \textbf{do not reset the
filter at the flip}. The half-difference measurement cannot tell $\pi$
apart anyway (that is the whole ambiguity), so the filter's near-zero
state is still valid --- it simply re-attaches to the true branch.
Resetting it by $\pi$ re-creates the shifted lock.
\srcline{Incident, diagnosis, fix: \code{teaching\_document.pdf} \S11; \code{--model micro --micro-pilots 9} reproduces it.}
\end{frame}

\begin{frame}{\proj{} Reimplementing \RN{}: does it reproduce?}
\textbf{Level 1 --- statistics level} (their model, faithfully:
Algorithms 1--2, MH weights, covariance mixing), at their setup
($N{=}20$, connectivity 0.2, $T{=}0.1$\,ms, 0\,dB):
\begin{center}
\small
\begin{tabular}{lcc}
\toprule
 & residual & \RN{, Eq.~27} bound\\
\midrule
DFPC & $8.0^\circ$ & $5.9^\circ$ (sparse-limit bound)\\
KF-DFPC & $2.1^\circ$ & ---\\
\bottomrule
\end{tabular}
\end{center}
Consensus converges to the average; the residual sits at the bound's
scale; filtering reduces it $\sim4\times$. \textbf{Both of the paper's
claims reproduce in its own model.}
\vspace{1ex}
\textbf{Level 2 --- physical level} ($N{=}2$, full simulator): with two
nodes the MH weight is $w_{12}=\tfrac12$, so ``consensus'' means
\emph{each node retunes its NCO by half its measured offset toward the
other}, both simultaneously --- now with real channels, latency, and
quantization. What happens is the next slide.
\srcline{\code{ota\_sync/dfpc.py}; tests in \code{tests/test\_ota\_sync.py}.}
\end{frame}

\begin{frame}{\proj{} Anti-phase capture, step by step}
\small
Naive OTA consensus: each node averages in its \emph{raw one-way}
measurement --- which contains $\phi_c$ (slide 3.6). In the relative
coordinate $r=\theta_A-\theta_B$, one round is
\[
r'=r+\tfrac12\wrapop(-r+\phi_c)-\tfrac12\wrapop(r+\phi_c).
\]
Walk through the default realization, $\phi_c=2.922$, $r_0=1.2$:
\[
\wrapop(1.2+2.92)=\wrapop(4.12)=-2.16
\qquad
\wrapop(-1.2+2.92)=+1.72
\]
\[
r_1 = 1.2+\tfrac{1.72}{2}-\tfrac{-2.16}{2}=1.2+0.86+1.08=3.14=\pi .
\]
One argument wrapped ($4.12>\pi$), the other did not --- their
disagreement is exactly $2\pi$, and half of it lands the pair at $\pi$.
And $\pi$ is \emph{stable}: there, both wraps give the same value
($-\pi\equiv\pi$), so $r'=\pi$ forever.
\begin{block}{Verified three independent ways}
full simulation ($r_1=-3.117$; deterministic across reruns) \ $\cdot$ \
a four-line iteration of the equation, no radio code \ $\cdot$ \ a seed
sweep matching the capture condition $|r_0|+|\phi_c|>\pi$ exactly
\end{block}
No wrap $\Rightarrow$ the halves cancel $\phi_c$ and the pair aligns:
the failure is \emph{realization-dependent}, and invisible in \RN{}'s
channel-free model.
\srcline{\code{--model dfpc}, \code{reciprocal=False}; verification: \code{teaching\_document.pdf} \S12.1.}
\end{frame}

\begin{frame}{\proj{} Head-to-head under identical physics}
\begin{center}
\small
\begin{tabular}{llccc}
\toprule
approach & origin & residual & gain & airtime\\
\midrule
two-tier micro-pilot & \proj{} & \textbf{28.1 mrad} & 99.98\% & 26.0\%\\
hybrid 1-way + anchors & \proj{} & 33.8 mrad & 99.97\% & \textbf{14.9\%}\\
two-way EKF & \proj{} & 80.8 mrad & 99.84\% & 19.1\%\\
KF-DFPC + reciprocity & \RN{} alg.\ + our PHY & 82.8 mrad & 99.83\% & 19.1\%\\
DFPC + reciprocity & \RN{} alg.\ + our PHY & 153.0 mrad & 99.42\% & 19.1\%\\
DFPC naive & \RN{} as published & 3009 mrad & 0.68\% & 19.1\%\\
\bottomrule
\end{tabular}
\end{center}
\begin{itemize}
\item ``+ reciprocity'' $=$ their update rule fed the channel-free
half-difference (using the side channel \RN{} already assumes).
\item \textbf{Unfiltered consensus is $2\times$ worse} because raw
frequency noise (0.442\,Hz/measurement) feeds the NCOs directly and
drifts $2\pi\!\cdot\!0.442\!\cdot\!0.05=139$\,mrad per interval;
$139\oplus45=146$ predicted, 153 observed.
\item \textbf{Filtered consensus $\approx$ our EKF} (82.8 vs.\ 80.8):
two different well-designed filtered loops reach the \emph{same}
physical floor --- the algorithm layer is interchangeable; the physics
of slide 4.3 binds.
\end{itemize}
\srcline{\code{--model compare}, seed 0.}
\end{frame}

\begin{frame}{\proj{} Design 1: two-tier micro-pilots --- the reasoning}
Chain of reasoning from the budget (slide on ablation):
\begin{enumerate}
\item Both dominant terms --- walk $kT$ and latency
$\sigma_{\hat\omega}T$ --- scale with the correction period $T$.
So correct more often.
\item But the full 4.6\,ms frame exists for \emph{acquisition}:
detection, timing, wide-range CFO. Once locked, timing drifts 0.016
samples per 10\,ms and frequency is known to $\sim$0.2\,Hz --- none of
that machinery is needed again.
\item Therefore: keep the full two-way frame once per 50\,ms, and
between frames send a \emph{bare} known sequence (ZC-255, 0.29\,ms) ---
correlate it at the predicted arrival, after derotating by the current
frequency estimate. Phase measured; corrections issued every
sub-interval.
\end{enumerate}
\begin{center}
\small
\begin{tabular}{lcccc}
\toprule
micro-pilots $M$ per interval & 0 & 2 & 4 & 9\\
\midrule
residual (mrad) & 71.5 & 35.3 & 27.9 & 22.8\\
airtime & 19.1\% & 22.6\% & 26.0\% & 34.6\%\\
\bottomrule
\end{tabular}
\end{center}
Scaling check: budget with $T\to T/(M{+}1)$ predicts 23--25\,mrad at
$M{=}4$; measured 27.9. \ $3\times$ tighter for $+7\%$ airtime.
\srcline{\code{ota\_sync/microsync.py}, \code{--model micro}. No counterpart in [RN23] (one pilot class; no airtime dimension).}
\end{frame}

\begin{frame}{\proj{} Design 2: hybrid calibration --- the reasoning}
Look again at slide 3.8: the \emph{only} thing the return direction
buys is separating $\theta$ from $\phi_c$. But the two behave
completely differently in time:
\begin{itemize}
\item $\theta$ wanders at the oscillator rate --- fast, power-law noise
\item $\phi_c$ between fixed stations is nearly constant
\end{itemize}
Paying double airtime every interval to re-measure a constant is waste.
So: \textbf{estimate $\phi_c$ as a state of its own}, with a 3-state EKF
over $[\theta,\omega,\phi_c]$:
\begin{itemize}
\item \textbf{One-way pilots} (cheap, every sub-interval) observe the
\emph{sum} $\theta+\phi_c$. Alone, they cannot split it --- the
direction $\theta{-}\phi_c$ is unobservable --- but the filter's priors
can: fast innovations get attributed to $\theta$ (large $q_\theta$),
slow ones to $\phi_c$ (small $q_c$).
\item \textbf{Two-way anchors} (expensive, every $K$ intervals) observe
$\theta$ and $\phi_c$ \emph{separately} (half-difference / half-sum)
and re-pin the split before prior-based attribution drifts.
\end{itemize}
The design question ``how often must we pay for reciprocity?'' becomes
``how fast does $\phi_c$ actually change?'' --- a \emph{channel}
question, not an oscillator question.
\srcline{\code{hybrid\_calibration/} package, \code{--model hybrid --anchor-every K}.}
\end{frame}

\begin{frame}{\proj{} How can priors split a sum the pilots cannot separate?}
\small
This deserves care, because it is the load-bearing idea of the hybrid
--- and its failure mode.
\begin{itemize}
\item A one-way pilot observes only $\theta+\phi_c$: its sensitivity to
the two states is \emph{identical} (Jacobian row $[1,\,0,\,1]$). From
one-way data alone, the direction $\theta-\phi_c$ is invisible.
\item So when a measurement disagrees with the prediction (an
\emph{innovation}), the filter cannot know which state moved. It
distributes the blame \textbf{in proportion to each state's prior
uncertainty}: the Kalman gain is built from $PH^\top$, so $\theta$'s
share $\approx P_{\theta\theta}$ and $\phi_c$'s share $\approx P_{cc}$.
\item Between anchors, $P_{\theta\theta}$ grows fast (oscillator
physics, large $q_\theta$) while $P_{cc}$ grows slowly (channel
believed static, small $q_c$) --- fast innovations get charged to the
clock, slow trends to the channel. \textbf{The prior is the separator.}
\item This predicts the failure mode before running it: if the channel
actually moves but $q_c$ says it cannot, the channel's share of blame
is $\approx0$; the clock absorbs everything, and the NCO chases the
channel. (The Doppler experiment two slides ahead shows exactly this.)
\item Anchors measure $\theta$ and $\phi_c$ \emph{separately},
collapsing both variances --- resetting the proportions before
prior-based attribution drifts too far.
\end{itemize}
\srcline{Standard Kalman-gain structure applied to the 3-state model of \code{hybrid\_calibration/hybrid.py}.}
\end{frame}

\begin{frame}{\proj{} Hybrid validation 1: static channel}
Prediction: with $\phi_c$ constant, anchors should barely matter.
\begin{center}
\begin{tabular}{cccc}
\toprule
anchor spacing $K$ & 1 & 5 & 20\\
\midrule
residual (mrad) & 31.8 & 32.9 & 34.0\\
airtime & 22.6\% & 14.9\% & \textbf{13.5\%}\\
\bottomrule
\end{tabular}
\end{center}
{\footnotesize (fully two-way control: 27.9\,mrad at 26.0\% airtime)}
\vspace{2ex}

A $20\times$ change in reciprocity cadence moves the residual by
2\,mrad, while sync airtime falls to 13.5\% --- below even the plain
two-way loop's 19.1\%. On a static link, reciprocity can be paid for
at 1\,Hz or slower. \textbf{Anchor cadence has decoupled from
oscillator quality.}
\srcline{Sweep in \code{hybrid\_calibration/README.md} (60 intervals, seed 0).}
\end{frame}

\begin{frame}{\proj{} Hybrid validation 2: moving channel (the converse)}
Motion changes the path length continuously, which shifts the received
frequency --- the \emph{Doppler} effect, $f_D=v/\lambda=3.05$\,Hz per
m/s at 915\,MHz. A frequency shift on the channel means its phase
$\phi_c$ now \emph{drifts}, up to $2\pi f_DT$ per interval --- so
anchors and priors \emph{should} start to matter:
\begin{center}
\small
\begin{tabular}{clcc}
\toprule
m/s & scheme & residual & gain\\
\midrule
0.2 & $K{=}5$, static prior & 446\,mrad & 95.3\%\\
0.2 & $K{=}5$, matched prior & 235\,mrad & 98.7\%\\
0.2 & $K{=}20$, static prior & 1752\,mrad & 55.1\%\\
0.5 & $K{=}1$, static prior & 128\,mrad & 99.6\%\\
0.5 & $K{=}1$, matched prior & \textbf{44\,mrad} & 99.95\%\\
\bottomrule
\end{tabular}
\end{center}
{\footnotesize (fully two-way controls: 43 / 33\,mrad --- immune, since
their measurements are channel-free)}
\vspace{1ex}

Why the static prior fails: the filter must decide whether a drifting
\emph{sum} means the clock moved or the channel moved. Told the channel
cannot move, it blames the clock --- and steers the NCO to chase the
channel. The matched prior ($\sigma_c\approx2\pi f_DT$) fixes the
attribution; sparse anchors are only partly rescued, because priors
re-attribute but only anchors truly \emph{separate}.
\srcline{Sweep in \code{hybrid\_calibration/README.md} (60 intervals, seed 0).}
\end{frame}

\begin{frame}{\proj{} Why user CSI absorbs static biases: two lines of algebra}
The claim from slide 3.7, now proved. At CSI time $t_0$ the user
measures each station's \emph{effective} channel --- true channel times
the station's carrier phase at that moment:
\[
g_i \;=\; h_i\,e^{j\phi_i(t_0)} .
\]
For joint transmission each station then weights its signal by
$g_i^{\,*}$ (transmit ``against'' the measured channel). What the user
receives later, at time $t$:
\[
\sum_i g_i^{\,*}\,h_i\,e^{j\phi_i(t)}
\;=\;\sum_i |h_i|^2\,e^{\,j[\phi_i(t)-\phi_i(t_0)]}.
\]
Everything static --- the channel phase, the one-way sync bias, chain
phases --- appears in both $\phi_i(t)$ and $\phi_i(t_0)$ and cancels
\emph{exactly}. What misaligns the sum is only the \textbf{drift since
the last CSI measurement}, and only its \emph{difference} between
stations (a common drift rotates everything together).
\vspace{1ex}

Measured with the one-way loop's residual drift as $\phi_i(t)$: mean JT
gain 100\% at 50\,ms CSI staleness, degrading only to 99.56\% at 1\,s.
One-way sync $+$ ordinary CSI is fully sufficient for this application
class.
\srcline{Evaluation: \code{--csi-gain} (\code{ota\_sync/coherent.py}).}
\end{frame}

\begin{frame}{\proj{} Which loop for which application}
Recall slide 3.7's split; now with numbers:
\vspace{1ex}
\begin{columns}[t]
\column{0.48\textwidth}
\textbf{CSI-aided (user feedback)}
\begin{itemize}
\item static biases absorbed by user CSI $\Rightarrow$ one-way suffices
\item what matters: differential drift between CSI updates
\item measured JT gain vs.\ CSI staleness: 100\% @ 50\,ms $\to$ 99.56\%
@ 1\,s
\end{itemize}
\column{0.48\textwidth}
\textbf{Open-loop (passive sensing)}
\begin{itemize}
\item carriers must align at the antennas $\Rightarrow$ two-way or
hybrid required
\item bounded closed-loop residual $\Rightarrow$ fixed 0.1--0.7\% gain
cost at \emph{any} integration time
\item dwell limited by target/channel dynamics, not by the stations'
clocks
\end{itemize}
\end{columns}
\vspace{2ex}
The second column is the property passive coherent sensing needs: an
\emph{uncorrected} pair of oscillators decorrelates and caps the dwell;
a closed loop holds the error bounded forever.
\srcline{JT evaluation: \code{--csi-gain} (\code{ota\_sync/coherent.py}).}
\end{frame}

\begin{frame}{Summary}
\begin{enumerate}
\item \textbf{Grading:} compare estimators to independent ground truth
(oracle path), never to their own inputs --- a $100\times$ difference
here \proj{}
\item \textbf{The floor:} walk $\sqrt{kT}$ $+$ latency
$\sigma_{\hat\omega}\ell T$ $+$ tracking; hand-predicted from physics,
then verified by ablation. The latency term is absent from
\RN{, Eq.~27} and unpublished through Dec.~2025 \proj{}
\item \textbf{The paper:} reproduces faithfully in its own model; over
a real channel its naive form is bistable (anti-phase capture --- walked
through and triple-verified), and its filtered form, given reciprocity,
meets our EKF at the same physical floor \proj{}
\item \textbf{Moving the floor:} cadence, not energy --- micro-pilots
(22.8--28\,mrad), and hybrid calibration (33\,mrad at 13--15\% airtime)
whose reciprocity cadence is set by channel coherence, validated under
Doppler in both directions \proj{}
\end{enumerate}
\vspace{1ex}
\textbf{Open:} channel-Doppler 4th state; $N{>}2$ wrap-safe consensus;
delayed-PLL literature check; \alert{hardware validation --- every
number in this deck is simulation}.
\end{frame}

\begin{frame}{References}
\footnotesize
\textbf{The reference paper}
\begin{itemize}
\item[{[RN23]}] M.~Rashid, J.~A.~Nanzer, ``Frequency and Phase
Synchronization in Distributed Antenna Arrays Based on Consensus
Averaging and Kalman Filtering,'' \emph{IEEE Trans.\ Wireless Commun.}
22(4):2789--2802, 2023. DOI 10.1109/TWC.2022.3213788.
\end{itemize}
\textbf{Classical results used in Part 1}
\begin{itemize}
\item D.~Chu, ``Polyphase codes with good periodic correlation
properties,'' \emph{IEEE Trans.\ Inf.\ Theory}, 1972.
\item T.~Schmidl, D.~Cox, ``Robust frequency and timing synchronization
for OFDM,'' \emph{IEEE Trans.\ Commun.}, 1997.
\item D.~Rife, R.~Boorstyn, ``Single-tone parameter estimation from
discrete-time observations,'' \emph{IEEE Trans.\ Inf.\ Theory}, 1974.
\item IEEE Std 1139 (frequency-stability metrics); J.~A.~Nanzer et al.,
\emph{IEEE TMTT} 65(5), 2017 (the $18^\circ$/90\% threshold).
\end{itemize}
\textbf{This project} (all in this repository)
\begin{itemize}
\item \code{teaching\_document.pdf} (full derivations, same order as
this deck); \code{simulation\_design.pdf};
\code{paper\_review\_rashid\_nanzer.pdf}; \code{LITERATURE\_REVIEW.md}
(verified 19-source prior-art pass); code: \code{ota\_sync/},
\code{hybrid\_calibration/}; entry point \code{simulation.py}; 18 tests.
\end{itemize}
\end{frame}

\end{document}
